%% shortname: Choice under Uncertainty, Lotteries and the Independence Axiom
\documentclass[11pt]{article}
\usepackage{microstyle}

\begin{document}

\lessonheader{15}{Chapter 6: choice under uncertainty --- lotteries, the simplex, continuity and the independence axiom}{October 26, 2026 (lab)}


\section{Midterm information}

\runin{Format.} 2 questions, 4 parts each:
\begin{flatnum}
  \item For a specific utility function, derive the Marshallian or Hicksian demand.
  \item A follow-up question regarding the solution.
  \item[3--4.] More follow-up, in more abstract settings.
\end{flatnum}

\runin{Advice.}
\begin{flatlist}
  \item Do not get stuck in part 3/4 --- part 4 can be answered without part 3.
  \item Explain answers in words.
  \item Review the assignments.
  \item Reading comprehension of technical words.
\end{flatlist}

\runin{Office hours.} Tuesday, 2:30--3:30.

\runin{Reading.} Chapter 6, Sections 6.A--6.D inclusive.


\section{Choice under uncertainty: lotteries}

First task: formulate a model that allows us to examine how an agent thinks about
and responds to uncertainty.

\topic{Outcomes}
Let $C$ denote a set of possible outcomes --- what happens once uncertainty has been
resolved. These could be consumption bundles, wealth, health, etc. Assume for now
that this set is finite, and let $N$ denote the number of outcomes,
$C=\{1,\dots,N\}$.

\topic{Simple lotteries}
A simple lottery is a list
\[ L=(p_1,p_2,\dots,p_N),\qquad p_n\ge 0,\qquad \sum_{n=1}^{N}p_n=1, \]
where $p_n$ denotes the probability that outcome $n$ occurs; they must sum to $1$ as
something must happen. The agent will have to choose between different lotteries,
each unique in the probabilities within it; using these choices we can recover her
preferences.

\topic{The simplex}
The feasible set of lotteries is
\[ \Delta=\Bigl\{p\in\R^{N}_{+}\ \Big|\ \sum_{n=1}^{N}p_n=1\Bigr\}. \]
Each lottery is a point in this $(N-1)$-dimensional simplex. The dimension is $N-1$
because the probabilities sum to $1$: if we know $N-1$ of them, we can calculate the
last one. The simplex is the feasible set of all lotteries over the possible
outcomes.

\topic{Degenerate lotteries}
A special lottery: a degenerate lottery has one of the probabilities equal to $1$,
meaning this outcome will happen for sure. It sits at a vertex (corner) of the
simplex.

\topic{Compound lotteries}
A compound lottery is a lottery whose outcomes are themselves lotteries. There could
be many lotteries in a row (a sequence), but there must be a finite end, eventually
in the set $C$. Compound lotteries can be reduced to simple lotteries.

\runin{Example.} $C=\{1,2,3\}$. With probability $\tfrac12$ the agent faces
$L_1=(\tfrac12,\tfrac12,0)$, with probability $\tfrac12$ she faces
$L_2=(\tfrac12,0,\tfrac12)$:
\[ p_1=\tfrac12\cdot\tfrac12+\tfrac12\cdot\tfrac12=\tfrac12,\qquad
   p_2=\tfrac12\cdot\tfrac12+\tfrac12\cdot 0=\tfrac14,\qquad
   p_3=\tfrac12\cdot 0+\tfrac12\cdot\tfrac12=\tfrac14, \]
\[ L=\bigl(\tfrac12,\tfrac14,\tfrac14\bigr). \]

\begin{center}
\begin{tikzpicture}
  % compound lottery tree
  \node[ecnode] at (-0.25,0) {$L$};
  \draw (0,0) -- node[ecnode,above left] {$\tfrac12$} (1.8,1.5);
  \draw (0,0) -- node[ecnode,below left] {$\tfrac12$} (1.8,-1.5);
  \foreach \y/\a/\b/\c in {1.5/\tfrac12/\tfrac12/0, -1.5/\tfrac12/0/\tfrac12} {
    \draw (1.8,\y) -- node[ecnode,above left,pos=0.7] {$\a$} (3.4,\y+0.9) node[ecnode,right] {$1$};
    \draw (1.8,\y) -- node[ecnode,above,pos=0.7] {$\b$} (3.4,\y) node[ecnode,right] {$2$};
    \draw (1.8,\y) -- node[ecnode,below left,pos=0.7] {$\c$} (3.4,\y-0.9) node[ecnode,right] {$3$};
  }
  % simplex, side 4.6: vertices 1, 2, 3; the distance to the side opposite vertex n is p_n (height = 1)
  \begin{scope}[xshift=5.4cm,yshift=-2cm,scale=1.15]
    \draw[budget] (0,0) node[ecnode,below left] {$1$} -- (4,0) node[ecnode,below right] {$2$}
                  -- (2,3.4641) node[ecnode,above] {$3$} -- cycle;
    \draw[dashed] (1.5,0.8660) -- node[ecnode,right] {$\tfrac14$} (1.5,0);
    \draw[dashed] (1.5,0.8660) -- node[ecnode,below left] {$\tfrac14$} (0.75,1.2990);
    \draw[dashed] (1.5,0.8660) -- node[ecnode,above left] {$\tfrac12$} (3,1.7321);
    \node[bundle] at (1.5,0.8660) {};
    \node[ecnode,above] at (1.5,0.9) {$L$};
  \end{scope}
\end{tikzpicture}
\end{center}
In the simplex (draw it with all sides equal, height $1$), the probability $p_n$ is
the perpendicular distance from $L$ to the side opposite vertex $n$.


\section{Preferences over lotteries}

Another axiomatic approach. Assume the agent evaluates compound lotteries as simple
lotteries in reduced form, as above: only the reduced lottery matters. Then we only
have to figure out how the consumer ranks simple lotteries.

Let $\mathcal{L}$ denote the set of (simple) lotteries and $\succsim$ the preference
relation on $\mathcal{L}$. Assume the agent is rational: preferences are complete and
transitive.

\topic{Continuity}
$\succsim$ is continuous if for any three lotteries $L,L',L''\in\mathcal{L}$ the sets
\[ \bigl\{\alpha\in[0,1]\ \big|\ \alpha L+(1-\alpha)L'\succsim L''\bigr\}
   \quad\text{and}\quad
   \bigl\{\alpha\in[0,1]\ \big|\ L''\succsim\alpha L+(1-\alpha)L'\bigr\} \]
are closed. Equivalently, consider the complement of the second set:
\[ \bigl\{\alpha\in[0,1]\ \big|\ \alpha L+(1-\alpha)L'\succ L''\bigr\}
   \ \text{is open}. \]
This means there is room for a small change in $\alpha$ and the (strict) preference
still holds: the ranking stays the same.

With these axioms (rational and continuous), $\succsim$ can be represented by a utility
function $U:\mathcal{L}\to\R$.

\topic{The independence axiom}
The next axiom leads to much more structure; it is a strong assumption and
controversial. Compare two lotteries $L$ and $L'$, and compare the compound lotteries
that mix each of them with the same third lottery $L''$. If $L\succsim L'$, then the
compound lotteries should follow suit.

$\succsim$ satisfies the independence axiom if for any $L,L',L''\in\mathcal{L}$ and
any $\alpha\in(0,1)$,
\[ L\succsim L' \quad\Longleftrightarrow\quad
   \alpha L+(1-\alpha)L''\ \succsim\ \alpha L'+(1-\alpha)L''. \]
That is, if $L\succsim L'$, then the mixture of $L$ and $L''$ is preferred to the
mixture of $L'$ and $L''$.

\begin{center}
\begin{tikzpicture}
  \draw (0,0) -- node[ecnode,above left] {$\alpha$} (1.4,0.8) node[ecnode,right] {$L$};
  \draw (0,0) -- node[ecnode,below left] {$1-\alpha$} (1.4,-0.8) node[ecnode,right] {$L''$};
  \node[ecnode] at (2.4,0) {$\succsim$};
  \draw (3.2,0) -- node[ecnode,above left] {$\alpha$} (4.6,0.8) node[ecnode,right] {$L'$};
  \draw (3.2,0) -- node[ecnode,below left] {$1-\alpha$} (4.6,-0.8) node[ecnode,right] {$L''$};
\end{tikzpicture}
\end{center}
The two branches are mutually exclusive: only one of them is realized. If the
$L''$ branch is realized, the agent gets $L''$ either way; so the comparison should
depend only on $L$ versus $L'$.


\section{Implications for the indifference curves}

\topic{1. Indifference curves are straight lines}
Assume $L\sim L'$ (i.e.\ $L\succsim L'$ and $L'\succsim L$). By independence this also
holds for the mixtures:
\[ \alpha L+(1-\alpha)L''\ \sim\ \alpha L'+(1-\alpha)L'', \]
and this should hold regardless of what $L''$ is.

Special case $L''=L$: then we conclude
\[ L\ \sim\ \alpha L'+(1-\alpha)L. \]
If $L\sim L'$, we should be indifferent about any mixture of the two (the betweenness
property). This means the indifference curve through $L$ and $L'$ contains the whole
segment between them: indifference curves are straight lines when $N=3$, and
(hyper)planes when $N\ge 4$.

\topic{2. Indifference curves are parallel}
Assume $L\sim L'$, fix $\alpha\in[0,1]$ and a third lottery $L''$, and let
\[ A=\alpha L+(1-\alpha)L'',\qquad B=\alpha L'+(1-\alpha)L''. \]
$A$ and $B$ lie at the same relative distance along the two sides $L''L$ and
$L''L'$, and $A\sim B$ by independence. The line connecting $A$ and $B$ must be
parallel to the line connecting $L$ and $L'$. So indifference curves shift in a
parallel way.

\begin{center}
\begin{tikzpicture}[scale=1.1]
  % left: simplex with parallel indifference lines
  \draw[budget] (0,0) -- (4,0) -- (2,3.4641) -- cycle;
  \begin{scope}
    \clip (0,0) -- (4,0) -- (2,3.4641) -- cycle;
    \foreach \c in {0.6,1.4,2.2,3.0,3.8} {
      \draw[thin] (\c-1,-0.6) -- (\c+3,3.4);
    }
  \end{scope}
  \node[ecnode] at (2,-0.45) {indifference lines in the simplex};
  % right: the triangle L, L', L'' with A, B
  \begin{scope}[xshift=5.8cm]
    \draw[budget] (0,0) node[ecnode,below left] {$L$} -- (3,0) node[ecnode,below right] {$L'$}
                  -- (2,3) node[ecnode,above] {$L''$} -- cycle;
    \draw[thick] (0,0) -- (3,0);
    \node[bundle] at (0.8,1.2) {}; \node[ecnode,left] at (0.75,1.2) {$A$};
    \node[bundle] at (2.6,1.2) {}; \node[ecnode,right] at (2.65,1.2) {$B$};
    \draw[dashed] (0.8,1.2) -- (2.6,1.2);
    \node[ecnode] at (1.5,-0.45) {$AB\parallel LL'$};
  \end{scope}
\end{tikzpicture}
\end{center}


\begin{transcriptnotes}
  \item The export is a single long page titled ``October 26 Lab'' (dated October 26,
        2026, 2:31 PM); it is filed as Lesson 15. Lesson 14 (normative representative
        agent, aggregation) has slides only, no handwritten export.
  \item Top of the page (midterm block): ``2:3:30 office hours -- Tuesday'' is read as
        Tuesday 2:30--3:30 (\textbf{unclear}: the start time). ``marshallian or
        hicksian demand'' kept; ``Explain answers means words'' $\to$ \emph{explain
        answers in words}; ``can answer part 4 without part 3'' merged with ``do not
        get stuck in part 3/4''. ``chapter 6, 6.A--6.D inclusive'' (MWG) set as the
        reading.
  \item ``ect.'' $\to$ \emph{etc.}; ``probabilty'' $\to$ \emph{probability};
        ``symplex'' $\to$ \emph{simplex}; ``compliment'' $\to$ \emph{complement};
        ``contraversal'' $\to$ \emph{controversial}; ``prefered'' $\to$
        \emph{preferred}; ``independent axiom'' $\to$ \emph{independence axiom};
        ``lotterys'' $\to$ \emph{lotteries}. ``his preferences'' typeset
        \emph{her preferences}, as in earlier lessons. The lottery is written
        $L(P_1,\dots,P_n)$ with $\sum P_i$; typeset $L=(p_1,\dots,p_N)$, $\sum_n p_n$,
        with $N$ the number of outcomes as defined on the page. $C=\{1,\dots,N\}$ is
        \textbf{added}.
  \item The simplex is written $A=\{p\in\R^N_+\mid\sum P_i=1\}$; typeset $\Delta$.
        \textbf{Corrected:} ``simplex = feasible set of possible outcomes'' --- the
        simplex is the set of \emph{lotteries} over the outcomes (the outcomes are its
        vertices). ``there must be a finite end'' read from ``must be a finite end
        eventually in the set of $C$''.
  \item Compound-lottery example: the hand-drawn tree has branch probabilities
        $\tfrac12,\tfrac12,0$ (top) and $\tfrac12,0,\tfrac12$ (bottom); the names
        $L_1$, $L_2$ and the computation of $p_3$ are \textbf{added} (the page gives
        $p_2=p_3$). The simplex sketch (``imagine all sides are equal'') is redrawn to
        scale; the red ``Bottom''/``Top'' labels on it are not reproduced. The
        explanation that $p_n$ is the distance to the side opposite vertex $n$ is
        \textbf{added}; checked: distances $\tfrac12,\tfrac14,\tfrac14$ for height 1.
  \item Continuity: the second set is written with $\precsim$; typeset with $L''$ on
        the left and $\succsim$, which is the same set. The pink margin note ``check
        continuity of the consumer on open \& closed sets'' is folded into the text.
        ``rank is same'' read as \emph{the ranking stays the same}. ``(rational and
        continuous)'' and $U:\mathcal{L}\to\R$ are \textbf{added}.
  \item Independence: the page states ``for any $L,L',L''$ and any $\alpha\in[0,1]$,
        if $L\succsim L'$ then \ldots''. \textbf{Corrected} to MWG Definition 6.B.5:
        $\alpha\in(0,1)$ and ``if and only if'' (at $\alpha=0$ both mixtures equal
        $L''$, so the ``only if'' direction would fail there). The page's statement is
        the ``if'' direction. The circled $L$, $L'$ and their ``$\vee$'' marks are the
        comparison of $L$ with $L'$; the sentence after the tree is expanded from
        ``two branches are mutually exclusive''.
  \item Implications: the page writes ``assume $L\sim L'$, $L\succsim L'$, $L'\succsim
        L$''; typeset ``$L\sim L'$ (i.e.\ \ldots)''. ``Between axiom'' $\to$
        \emph{betweenness property}. ``line when $n=3$, hyperplane when $n=4$'' is
        typeset with $N$ and ``$N\ge 4$''. \textbf{Corrected:} $B=\alpha L'+(1-\alpha)L'$
        has $L'$ twice; it is $\alpha L'+(1-\alpha)L''$. ``$(-\alpha)$'' $\to$
        $(1-\alpha)$. ``same distance to the two sides'' typeset as the same
        \emph{relative} distance along the two sides.
  \item Figures: the left sketch (a simplex with $L$, $L'$, $L''$ and two angle marks
        for similar triangles) is redrawn as a simplex with parallel indifference
        lines; the right sketch ($L$, $L'$, $L''$ with $A$, $B$) is redrawn to scale
        with $\alpha=0.6$.
\end{transcriptnotes}

\end{document}
